\section{Introduction}

Single-pixel (ghost) imaging \cite{Pittman1995PRA,Bennink2002PRL,Shapiro2008PRA,Bromberg2009PRA} has become the method of choice for imaging through media that defeat conventional cameras \cite{Bina2013PRL,Cheng2009OE}: turbid water, fog, dynamic scattering media \cite{PengChen2024APL}, complex channels around corners \cite{PengXiaoChen2025OE}. The instrument is disarmingly simple --- structured illumination patterns and a bucket detector \cite{Duarte2008SPM,Edgar2019NatPhoton} --- and the textbook design instinct is equally simple: use an ordered orthogonal basis. A Hadamard scan acquires one transform coefficient per frame \cite{Zhang2017OE,Yu2019Sensors}, the measurement matrix is perfectly conditioned, and an exact inverse waits at the end of the acquisition; against additive noise this is provably the safe choice.

Dynamic media break this instinct in a specific, instructive way. A time-varying channel multiplies the bucket record by a slowly drifting gain: the $n$-th measurement is $R_n = a_n B_n$, with $B_n$ the object-dependent bucket coefficient and $a_n > 0$ an unknown per-frame gain whose log-gain $\ell_n = \log a_n$ is confined to a slow-drift class $S$ of dimension $p$. Ghost imaging through dynamic and complex scattering media \cite{PengChen2024APL} shows a self-calibrating scheme built on \emph{random} illumination with supervised correction of dynamic scaling factors succeeding where ordered scans fail, and optical wireless transmission around corners \cite{PengXiaoChen2025OE} relies on the same random-pattern robustness. Read naively, this is a puzzle: the ordered Hadamard scan is the better-conditioned instrument, yet randomness --- usually a concession, not a feature --- wins; it is desirable to determine when, and by what mechanism, a drifting gain becomes blindly separable at all.

The resolution is that under gain drift the ordering itself becomes the liability. An ordered scan sorts the object's transform coefficients into a smooth sequence in acquisition time, so the object acquires a \emph{temporal signature}: a slow gain trace and a slow coefficient trace cast the same shadow on the one-dimensional bucket record, and no estimator can tell them apart from that record alone. Hadamard is not defeated by randomness; it is defeated by chronology. This confounding is not an artifact of any particular correction algorithm but a genuine algebraic ambiguity --- for a square unconstrained design, \emph{any} invertible pattern matrix can absorb a nonconstant gain into a perfectly valid alternative object (Theorem~A). Conversely, random carriers escape it statistically when their log-carrier window means are time-homogeneous and concentrate as condition ($\star$): slow gain is then the only persistent local modulation, and relative gain is recovered up to the unavoidable global scale (Theorem~B). Self-calibration is thus a property one can \emph{design into acquisition time}, with price, rates, and failure boundaries computable in advance.

We develop that observation into a quantitative identifiability theory; the contributions are:

\begin{enumerate}[leftmargin=*,itemsep=1pt,topsep=2pt]
\item \textbf{A three-notion identifiability framework} for the bilinear model $R_n = a_n B_n$, separating exact-algebraic identifiability, statistical (asymptotic) identifiability, and estimation conditioning. Within it, the square-design obstruction (Theorem~A) and the tall-design thresholds (Theorem~A$'$) are proved for every fixed drift class $S$ containing the constants: generic exact identifiability at $N \ge K + p$ for almost every design--object pair, uniform identifiability over every nonvanishing-carrier object at $N \ge 2K + p - 1$, and a generic local threshold at $N = K + p - 1$ that is an if-and-only-if for $K \ge 2$, with genuine positive-dimensional alias families below it (Secs.~4, 9). This is distinct from generic bilinear-injectivity results \cite{KechKrahmer1603.07316,LiLeeBresler1501.06120,AhmedRechtRomberg1211.5608,ChoudharyMitra1402.2637,KuengRauhutTerstiege1410.6913}: the operator here is a diagonal-row map, and the anchor is temporal, not structural.

\item \textbf{The stationarity anchor (Theorem~B)} and its consistency condition ($\star$): under explicit triangular-array bias and window-variance conditions, random carriers pin the relative gain statistically. A separate \emph{oracle-known-carrier} calibrated soft-log theorem covers the photon-limited case, with a conditional $1/(W\bar\lambda)$ law in the locally-flat-gain regime. The blind high-count and known-carrier low-photon experiments are distinct.

\item \textbf{A synthesis, not a trade-off:} randomized Hadamard (SRHT) designs retain orthogonal conditioning while randomizing the temporal carrier. Pixel signs and a pixel permutation control ensemble covariance through an exact Walsh-domain criterion; a separate random acquisition order supplies finite-population exchangeability. A conditioned draw is deterministic, so these two randomization semantics are not conflated.

\item \textbf{An exact finite-noise reconstruction bridge} (Theorem~1) expressing raw relative MSE for any fixed linear reconstructor through a leverage functional $B_L$, under explicit exogenous-gain hypotheses. Physical nonnegative Poisson measurements, signed coefficient-domain noise, and complementary-pair noise are kept as distinct specializations \cite{Ferri2010PRL,Sun2012OE}.

\item \textbf{A finite-$N$ flip diagnostic} $\rho^*$ in $\rho_{\mathrm{pair}}$ (Prop.~3): its gain-known true-$S_1$ skeleton is validated without fitted parameters; the production normalization and same-record blind arm are reported separately as robustness and conditional-proxy checks.
\end{enumerate}

All results are theoretical, supported by numerical evidence; no hardware validation is reported. Identifiability near $N = K + p$ is a uniqueness statement, not a conditioning guarantee, and random square designs remain algebraically non-identifiable --- their advantage is statistical. The organizing idea is the distilled paradox: one makes the gain observable not by measuring it directly, but by making the object stop pretending to be time.

Figure~1 maps ordered-Hadamard, iid-random, tall-random, and SRHT designs on the resulting conditioning-versus-identifiability axes; Sec.~10 collects the design rules, scope, and outlook.

\begin{figure}[!t]
\centering
\includegraphics[width=0.55\columnwidth]{fig1_conceptmap.png}
\caption{Two-axis map of inversion conditioning against blind identifiability. Ordered orthogonal scans (Hadamard/DCT/Fourier) are well-conditioned yet blindly non-identifiable; i.i.d.-random (DGI) acquisition buys the statistical anchor at the cost of conditioning; a single fixed pattern (grey) has neither; tall random designs ($N\ge K+p$) restore exact identifiability modulo scale for almost every design--object pair (Theorem~A$'$(ii)); SRHT (star) retains orthogonal conditioning while randomizing the temporal carrier. The dashed diagonal marks the apparent trade-off the SRHT synthesis dissolves (Secs.~4 and 5).}
\end{figure}

\section{Problem setup and three notions of identifiability}

In a single-pixel (bucket-detection) acquisition of $N$ frames, frame $n$ illuminates a static object $T \in \mathbb{R}^K$ with a programmed pattern $m_n \in \mathbb{R}^K$, producing the ideal bucket value $B_n = m_n^{\mathsf T} T$ --- the \emph{object carrier}. The record is corrupted by a slowly drifting multiplicative gain (source power, detector responsivity, coupling):

\[
R_n \;=\; a_n B_n, \qquad a_n = e^{\ell_n} > 0 ,
\]

where $a_n$ is the unknown per-frame gain and the log-gain trajectory $\ell = (\ell_1,\dots,\ell_N)$ is confined to a \emph{drift class} $S \subset \mathbb{R}^N$, $p = \dim S$, always containing the constants. The constant direction is a pure gauge: $(a, T) \mapsto (c\,a,\; T/c)$ leaves every $R_n$ invariant, so all statements are up to a global scale. For physically slow drift, $S$ is a low-pass (Fourier) space of normalized bandwidth $\rho_{\mathrm{bw}}$, $p \approx \rho_{\mathrm{bw}} N$ (not to be confused with the adjacent-pair decorrelation $\rho_{\mathrm{pair}}$ of Sec.~8). The blind problem --- recover $(a, T)$ from $\{R_n\}$ alone --- is bilinear, and asking whether it is ``solvable'' turns out to be three distinct questions; separating them is itself a contribution of this paper.

\begin{nameddef}{Conventions (fixed throughout)}
$K$ --- pixels; $T\in\mathbb{R}^K$ the object, nonnegative unless stated; $N$ --- frames; $S$ --- linear log-gain drift class, \textbf{always containing the constants}, $p=\dim S$. \textbf{All} identifiability statements are modulo the gauge $(a,T)\mapsto(c\,a,\,T/c)$, fixed by centering the log-gain. Theorem~1 and Prop.~3 use raw relative MSE; least-squares scale-aligned MSE is a separately labeled image-quality metric. \emph{Local-differential}, \emph{generic-exact}, and \emph{uniform-exact} identifiability are distinct properties with distinct thresholds (Theorem~A$'$); none implies conditioning. $\rho_{\mathrm{bw}}$ --- normalized gain bandwidth, $p\approx\rho_{\mathrm{bw}}N$ (Sec.~4); $\rho_{\mathrm{pair}}$ --- adjacent-pair gain decorrelation (Sec.~8); related but not equal. The object functionals $K_{\mathrm{eff}}$, $K_\infty$, $K_4$ and pipeline constants $C_0$, $D_H$, $B_L$ are defined at first use (Prop.~1, Secs.~6--8); a consolidated notation table appears in Supplementary Material S5.
\end{nameddef}

Pattern entries $m_n$ may be signed (Hadamard) or nonnegative (speckle); where signed coefficients appear as carriers, the Sec.~5 bucket conventions specify their physical realization.

\noindent\textbf{Algebraic identifiability.} Does the noiseless record determine $(a,T)$ up to gauge, or does some other pair with $\ell' - \ell \in S$ nonconstant reproduce the data exactly --- a finite-sample, deterministic question about collisions $MT = D_s M T'$? Section~4 shows that every square invertible design, however well-conditioned, fails it, and that overdetermination restores it generically at $N \ge K + p$.

\noindent\textbf{Statistical identifiability.} Under random pattern ensembles a weaker but operationally decisive property becomes available: if the carrier is time-homogeneous and its log-carrier window means concentrate according to ($\star$), then any persistent slow modulation of those averages is attributable to gain, and the relative gain is recoverable even where exact algebraic separation is impossible. Stationarity alone is not sufficient without the window-variance condition. This stationarity anchor (Sec.~5, Theorem~B) is the mechanism the field's empirical successes have been implicitly exploiting.

\noindent\textbf{Estimation conditioning.} Uniqueness is not accuracy: even when one of the first two notions holds, the finite-noise error of a concrete estimator is governed by conditioning quantities --- the window $W$, the residual-gain variance $v$, the leverage $B_L$ --- quantified in Secs.~6 and 8. Identifiability at the threshold $N = K + p$ says nothing about stability there.

Much of the confusion surrounding blind gain correction stems from conflating these three notions: judging a scheme by the visual quality of a blindly corrected reconstruction entangles algebraic ambiguity, the statistical anchor, and estimator conditioning into a single scalar --- three failures with different remedies. Figure~1 disentangles them as a two-axis map --- inversion conditioning against blind identifiability --- on which ordered Hadamard, i.i.d.-random (DGI), tall random, and SRHT acquisitions occupy visibly different corners (Secs.~4 and 5).

\section{Relation to prior work}

Gain--object separation is a bilinear inverse problem: lifting the unknowns to the rank-one array $X=aT^{\top}$, each bucket $R_n=a_n B_n$ with $B_n=m_n^{\top}T$ is a bilinear form in the object $T$ and the gain trace $a$. This places the question inside a mature literature on bilinear identifiability; we state what it settles and what it leaves open for the diagonal-gain geometry.

The sharpest injectivity counts for lifted bilinear maps are Kech--Krahmer's optimal conditions for generic identifiability \cite{KechKrahmer1603.07316}, requiring on the order of $2(n_1+n_2)-4$ generic rank-one measurements. These counts do not transfer here: the operator $L_M(X)_n=a_n\,m_n^{\top}T$ is a \emph{diagonal-row} map, not a generic rank-one measurement, and the relevant question is the intersection geometry of the object subspace with its gain-modulated copy, $\,U\cap D_sU$. Choudhary--Mitra's scaling laws \cite{ChoudharyMitra1402.2637} are the closest in spirit to our Theorem~A: their rank-two nullspace obstruction is exactly the mechanism by which a square unconstrained design absorbs a nonconstant gain into a redefined object. Li--Lee--Bresler's unified framework \cite{LiLeeBresler1501.06120} supplies the subspace- and sparsity-constrained thresholds that our tall-design counts ($N\ge K+p$ generic, $N\ge 2K+p-1$ uniform) specialize to for the diagonal-gain model. Ahmed--Recht--Romberg \cite{AhmedRechtRomberg1211.5608} establish convex recovery for blind deconvolution under subspace priors, and Kueng--Rauhut--Terstiege \cite{KuengRauhutTerstiege1410.6913} give low-rank-from-rank-one-measurement guarantees; both are the algorithmic backdrop for our purely algebraic and statistical criteria.

On the imaging side, single-pixel practice has long weighed ordered Hadamard against random speckle for conditioning and noise, but without an identifiability account of slow multiplicative drift. Learned correction of dynamic scaling factors \cite{PengChen2024APL} and learned encoding for optical wireless links \cite{PengXiaoChen2025OE} address the same nuisance empirically; our analysis explains why their \emph{static-correction ceiling} --- the error floor left by any time-independent correction applied to a drifting record --- is the honest bound, and when a drifting gain becomes blindly separable at all. Blind self-calibration of sensor gains and phases has a long history outside imaging --- eigenstructure methods \cite{WeissFriedlander1990}, blind sensor-network calibration \cite{Balzano2007IPSN}, biconvex compressive self-calibration \cite{LingStrohmer2015} --- and Balzano--Nowak similarly find gains recoverable under mild oversampling, which our threshold $N\ge K+p$ makes exact for the diagonal bilinear model. On the acquisition-design side, pink-noise speckle \cite{Nie2021PRA} and sparse Hadamard sequencing \cite{Han2025OE} optimize the pattern set directly but target additive-noise robustness, not multiplicative-drift identifiability.

Against this background the contribution is specific: we do not add to generic bilinear injectivity, settled above. New is the temporal-statistical mechanism --- relative-gain identifiability secured by making the carrier stationary (Theorem~B), a general mixing-law upper bound together with a matching fixed-iid-Gaussian interior benchmark, the ordered-orthogonal confounding failure mode it repairs, the SRHT synthesis, and the finite-noise bridge that renders the phase diagram quantitative. The novelty is the stationarity anchor, not bilinear injectivity per se.

\section{The algebraic obstruction and when oversampling cures it}

The paradox of Sec.~1 is not an estimator artifact. Stack $N=K$ frames of the noiseless record $R_n = a_n B_n$ into an invertible square design $M$: the multiplicative drift is absorbed \emph{exactly} into a relabelled object.

\begin{namedthm}{Theorem A (square non-identifiability)}
For any invertible $M\in\mathbb R^{K\times K}$ and any gain $a=(a_n)$ nonconstant on the support of the carrier $MT$, there exists a distinct object $T'\neq cT$ producing the identical bucket record: $\operatorname{diag}(a)MT = M T'$ with $T'=M^{-1}\operatorname{diag}(a)MT$. Gain--object separation is unidentifiable up to the global-scale gauge $a\mapsto ca,\ T\mapsto T/c$. (Proof: Appendix~A.)
\end{namedthm}

The object is a barcode printed onto time and the drift a smudge across it; when the barcode carries slow structure, the smudge can be \emph{reread as a different barcode} rather than dismissed as blur --- a genuine orbit of the bilinear map, for \emph{every} invertible design. Two corollaries sharpen the trap (full statements and proofs: Appendix~A). \emph{Corollary~1 (slowness does not help):} restricting $\ell\in S$ to a slow drift class does not restore identifiability at $N=K$ --- for any nonconstant $s\in S$ the alternative pair $\big(\ell-s,\; T'=M^{-1}\operatorname{diag}(e^{s})MT\big)$ explains the record exactly. \emph{Corollary~2 (known support zeros):} nonnegativity alone does not break the orbit. For fixed known support, a displayed quotient-Jacobian full-rank condition is equivalent to \emph{differential} identifiability and is sufficient for ordinary local identifiability; the exact ordinary-local condition is isolation of the nonlinear gauge-fixed support zero. The count $K_0\ge p-1$ is necessary only for that full-rank test and is not by itself an ordinary-identifiability theorem. This qualification does not affect Theorem~A$'$(i) below.

The first escape is overdetermination. For a tall design $M\in\mathbb R^{N\times K}$, $N>K$, a collision $\operatorname{diag}(e^{s'})MT'=\operatorname{diag}(e^{s})MT$ pits $K+p$ unknowns against $N$ equations, suggesting uniqueness once $N\ge K+p$; the following theorem turns this count into proved sufficiency thresholds, for every fixed drift class $S\ni\mathbf 1$, at the price of almost-everywhere quantifiers with explicit Lebesgue-null exceptional sets.

\begin{namedthm}{Theorem A$'$ (proved tall-design identifiability results)}
Let $K\ge1$, $N\ge K$, and let $S\le\mathbb R^N$ be a fixed, known linear subspace of dimension $p\ge2$ with $\mathbf 1\in S$. For $M\in\mathbb R^{N\times K}$ define $\Phi_M(T,s)=D_{e^{s}}MT$ and $\mathcal U_M=\{T\in\mathbb R^K:(MT)_n\neq0\ \text{for every } n\}$. Parameters are identified modulo the positive global-scale gauge $(T,s)\sim(cT,\,s-\log c\,\mathbf 1)$, $c>0$.

\emph{(i) (Generic local threshold and genuine failure below it, $K\ge2$.)} There is a Lebesgue-null set $E_{\mathrm{loc}}\subset\mathbb R^{N\times K}\times\mathbb R^K$ such that for every $(M,T)\notin E_{\mathrm{loc}}$ with $T\in\mathcal U_M$ and every $s\in S$: $\operatorname{rank}D\Phi_M(T,s)=\min\{N,K+p-1\}$. Consequently $(T,s)$ is locally identifiable modulo gauge iff $N\ge K+p-1$. If $N\le K+p-2$ then a gauge-normalized neighborhood of $(T,s)$ contains a real-analytic manifold of exact aliases of dimension $K+p-N-1\ge1$, and every other sufficiently nearby point of it is gauge-inequivalent.

\emph{(ii) (Generic exact-identifiability sufficiency.)} If $N\ge K+p$, there is a Lebesgue-null set $E_{\mathrm{gen}}\subset\mathbb R^{N\times K}\times\mathbb R^K$ such that for every $(M,T)\notin E_{\mathrm{gen}}$ and every $s\in S$: $D_{e^{s'}}MT'=D_{e^{s}}MT$ with $(T',s')\in\mathbb R^K\times S$ implies $T'=cT$, $s'=s-\log c\,\mathbf 1$ for a unique $c>0$. Exact identifiability holds simultaneously for every possible true drift $s\in S$. This statement is \textbf{not} uniform over all objects for one fixed $M$.

\emph{(iii) (Uniform exact-identifiability sufficiency on the nonvanishing-carrier set.)} If $N\ge2K+p-1$, there is a Lebesgue-null set $E_{\mathrm{unif}}\subset\mathbb R^{N\times K}$ such that for every $M\notin E_{\mathrm{unif}}$, every $T\in\mathcal U_M$, every $s\in S$, and every $(T',s')\in\mathbb R^K\times S$: $D_{e^{s'}}MT'=D_{e^{s}}MT$ implies $T'=cT$, $s'=s-\log c\,\mathbf 1$ for a unique $c>0$. The same generic design works simultaneously for all nonvanishing-carrier objects and all true drifts in $S$.

\emph{(iv) (Degenerate one-dimensional object stratum.)} If $K=1$ then for every $M$ and every $T\in\mathcal U_M$ and every $s\in S$, the pair $(T,s)$ is globally identifiable modulo gauge, for every admissible $N$. Neither $N\ge K+p$ nor $N\ge2K+p-1$ is necessary when $K=1$.

(Proof: Appendix~B.)
\end{namedthm}

\noindent\textbf{Sharpness status.} Parts (ii) and (iii) are sufficient conditions, not proved if-and-only-if thresholds. For $K\ge2$, the below-wall part of (i) proves generic exact failure whenever $N\le K+p-2$. The remaining necessity assertion for the $K+p$ bound is the boundary case $N=K+p-1$: dimension counting predicts nonlocal aliases there, but dominance of the bad-pair projection has not been proved. Likewise, dimension counting predicts that a generic design fails uniform identifiability for some object when $N\le2K+p-2$, but this necessity direction has not been proved in the intermediate regime. Thus $K+p$ and $2K+p-1$ are conjecturally sharp for $K\ge2$ under generic/nondegenerate drift geometry, not universal sharp thresholds for every fixed $S$. The $K=1$ stratum explicitly rules out any universal necessity claim.

The proof is a parametric-transversality argument (Appendix~B): a rowwise perturbation of $M$ acts as a global witness on a collision incidence manifold, and no genericity of the drift family is invoked, so the statements hold verbatim for the physical Fourier low-pass class $S_{\mathrm{LP}}$ and any other fixed drift convention containing the constants. Empirically, the $K=64$ and $K=128$ threshold scans of Sec.~9 (Fig.~6) show an exact $0\%\to100\%$ rank transition at $N=K+p-1$ across all five drift dimensions at both object sizes (3250 cells, no exceptions) --- a direct verification of the generic local rank wall of part (i). Part (ii) asserts \emph{generic exact finite-sample identifiability}, not a conditioning bound and not an algorithmic recovery guarantee; stability near the thresholds is the separate third notion of Sec.~2, quantified in Secs.~6 and 8.

The picture is a dichotomy. When $N\ge K+p$, overdetermination \emph{alone} separates gain from object for almost every design--object pair --- no statistical mechanism is logically needed, though randomization still buys stable conditioning (Fig.~1). When $K<N\le K+p-2$, generic local and exact identifiability fail through positive-dimensional nearby alias families; at $N=K+p-1$, local identifiability holds generically, while failure of global exact identifiability remains a dimension-count prediction. In that under-resourced window some additional information is required --- the stationarity anchor of Sec.~5 and support/sparsity priors are examples, not an exhaustive list. The regimes meet through the drift bandwidth: a gain occupying a fraction $\rho_{\mathrm{bw}}$ of the frame band spans $p\approx \rho_{\mathrm{bw}}N$ low-pass modes, so $N\ge K+p$ reads as $N \gtrsim K/(1-\rho_{\mathrm{bw}})$. Slow drift ($\rho_{\mathrm{bw}}\ll1$) is cured by modest oversampling; as $\rho_{\mathrm{bw}}\to1$, $p\to N$ and blind algebraic separation is impossible for \emph{any} $N$, forcing some non-algebraic information source.

\section{Statistical identifiability: the stationarity anchor}

Section~4 bought identifiability with surplus rows; this section develops a second, logically distinct escape that operates even at fixed $N$ and explains \emph{why} randomized acquisition self-calibrates while ordered acquisition does not. A drifting gain is a boat carried by a slow current, and the bucket record is our only view of the water: if the waves change character with position --- as the ordered Hadamard coefficient sequence does --- we cannot tell the current from the coastline; if the waves are \emph{statistically identical everywhere}, any slow modulation of the record has exactly one explanation. Local averaging then pins the \emph{relative} drift, the absolute scale remaining gauged away. This intuition, and nothing stronger, is what the results below formalize.

The log-domain record can be described by $Y_n=\log R_n=\ell_n+\log B_n$. The anchor requirement is a condition on the carrier alone:

\begin{nameddef}{Condition ($\star$) (stationary carrier, quantitative form)}
There is a scalar $m_T$ --- depending on the object but not on time, and absorbed by the global-scale gauge --- such that the windowed carrier averages over windows $W_n$ of length $W$ satisfy
\[
\begin{gathered}
\max_n \Big| W^{-1}\!\!\sum_{k\in W_n}\!\log B_k \;-\; m_T \Big| \;\longrightarrow\; 0\\
\text{in probability as } W\to\infty,\ W/N\to 0 .
\end{gathered}
\]
Condition ($\star$) and Theorem~B presuppose a positive carrier with an integrable log on the analyzed record, typically enforced by a physical positivity/offset margin. Photon-counting records with zeros are treated separately by the oracle-known-carrier soft-log Theorem~C; that result is not a blind extension of ($\star$).
\end{nameddef}

The criterion is deliberately \emph{uniform} (a maximum over all window locations); the pointwise-to-uniform upgrade typically costs a $\log N$ refinement (the $\delta\mapsto\delta/N$ step of Appendix~C).

I.i.d. random illumination satisfies ($\star$) under the stated log-tail control and a window sequence large enough to pay the $\log N$ uniformity cost; iid alone gives only the pointwise law without that scaling. A uniform acquisition-order permutation $P_{\mathrm{row}}$ makes the sequence finite-population exchangeable over that randomization draw, whereas a shared pixel permutation $P_{\mathrm{col}}$ controls ensemble covariance but is not joint row exchangeability. After conditioning on either draw the record is deterministic, and a particular realization can fail the window criterion. For the i.i.d. case:

\begin{namedthm}{Proposition 1 (random-basis carrier)}
For i.i.d. random patterns with pixel mean $\mu_I$ and standard deviation $\sigma_I$, the carrier $\{B_n\}$ is white and strictly stationary, and the object enters its first two moments only through the scalars $S_1=\sum_j T_j$ and $S_2=\sum_j T_j^2$: $\mathbb{E}B_n=\mu_I S_1$, $\operatorname{Var}B_n=\sigma_I^2 S_2$, hence $\mathrm{CV}_B=(\sigma_I/\mu_I)/\sqrt{K_{\mathrm{eff}}}$ with $K_{\mathrm{eff}}=S_1^2/S_2$. The Gaussian approximation additionally requires the Lindeberg-type spikiness condition $K_\infty=S_2/\|T\|_\infty^2\to\infty$; the log-domain expansion further requires positivity, log-integrability, and the lower-tail control stated in Appendix~C. The reduction of the object to two scalars holds at the leading Gaussian/log-variance level, not exactly. Higher cumulants can depend on object shape beyond $(S_1,S_2)$: their contribution to the log mean is a time-constant gauge term, while their contribution to log variance and tails changes finite-sample rate constants. It does not create a time-varying carrier bias or invalidate consistency when ($\star$) holds. (Proof: Appendix~C.)
\end{namedthm}

Figure~2 shows the two carriers side by side: the random trace is a time-homogeneous, object-dependent noise band without an object-specific temporal envelope, whereas the ordered Hadamard trace \emph{is} the object's coefficient envelope --- the object masquerades as a temporal trend, precisely the degree of freedom a slow gain occupies, and ($\star$) fails. The condition's logical status:

\begin{namedthm}{Proposition 2 (($\star$) is an estimator criterion, not a universal one)}
Condition ($\star$) is necessary and sufficient for consistency of the stated blind windowed log estimator. It is \emph{not} an if-and-only-if condition for identifiability by all mechanisms or for every windowed correction rule: the overdetermined designs of Sec.~4 identify algebraically --- for almost every design--object pair at $N\ge K+p$, Theorem~A$'$(ii) --- without any stationarity, and conversely ($\star$) confers no exact algebraic identifiability. (Proof: Appendix~C.)
\end{namedthm}

\begin{figure*}[!t]
\centering
\includegraphics[width=0.50\textwidth]{fig2_carrier_traces.png}
\caption{Carrier stationarity: noiseless bucket carriers $B_n$ (thin), running means (bold), and the $\pm1$\,SD envelope for one object (digit\_0; 8192 frames, $64\times64$), before any gain is applied. (a) The i.i.d.\ random-uniform carrier is a time-homogeneous, object-dependent noise band without an object-specific temporal envelope; (b) the ordered Hadamard trace \emph{is} the object's coefficient envelope --- a deterministic, non-stationary temporal signature. Rejection counts and the Parseval invariant: Sec.~9.1.}
\end{figure*}

\noindent\textbf{Bucket conventions.} Three bucket-signal conventions coexist (full statement: Supplementary Material S5): (i) \emph{signed coefficients} --- $B_n$ can vanish or change sign, so $\log B_n$ is undefined; (ii) \emph{physical positive offset} --- $B_n\ge B_{\min}>0$ with an integrable log; Theorem~B and its Sec.~6 windowed-log estimator live here under (B1)--(B4); (iii) \emph{$\pm$ pairwise} --- complementary pairs differenced, drift entering through $\rho_{\mathrm{pair}}$ (Sec.~8). The raw signed arms of Sec.~9 follow (i) and are \emph{masked diagnostics only}, not tests of Theorem~B.

\begin{namedthm}{Theorem B (statistical relative-gain consistency)}
For the triangular array $Y_{N,n}=\ell_{N,n}+m_{T,N}+z_{N,n}$, let $b_{N,n}(W)$ be the actual deterministic smoothing bias and $V_{N,n}(W)$ the actual variance of the rowwise-stationary carrier window. If $W_N\to\infty$, $W_N/N\to0$, $b_{N,n}(W_N)\to0$, and $V_{N,n}(W_N)\to0$, then the windowed estimator is pointwise consistent, with
\[
\mathbb E[\hat g_{N,n}-(\ell_{N,n}+m_{T,N})]^2=b_{N,n}(W_N)^2+V_{N,n}(W_N).
\]
For normalized-time H\"older radius $L_N$ and absolutely summable rowwise covariance, sufficient conditions are $L_N(W_N/N)^\beta\to0$ and $\sigma_{\mathrm{abs},N}^2/W_N\to0$. Centered-gain recovery requires these bounds uniformly over the indices used in centering. For $\beta>1$, the claimed bias order applies only to centered/moment-cancelling interior weights or a defined boundary local-polynomial rule. The absolute scale remains unidentifiable. (Proof: Appendix~C.)
\end{namedthm}

Theorem~B must be read with its scope intact: it is an asymptotic, high-probability statement about \emph{relative}-gain recovery up to the gauge --- never exact finite-sample identifiability of $(a,T)$ --- so it does not contradict Theorem~A: a random square design remains algebraically non-identifiable; randomness buys statistical separability, a different currency. This also distinguishes the mechanism from the bilinear-inverse literature \cite{LiLeeBresler1501.06120,KechKrahmer1603.07316,ChoudharyMitra1402.2637}: the anchor is a property of the carrier's \emph{temporal distribution}, not of subspace injectivity.

The anchor also explains why the self-calibrating scaling-factor correction of Ref.~\cite{PengChen2024APL} can work at all: under Proposition~1's Gaussian approximation the bucket record is a slow trend plus stationary residuals, the same decomposition that motivates an SCGI-style Gaussian-likelihood fit. Theorem~B proves consistency only for the stated windowed log estimator; consistency of the likelihood fit would require a separate M-estimation argument.

Figure~3 is the core numerical evidence: under random illumination the blind gain error $\|\hat a-a\|/\|a\|$ tracks $\mathrm{CV}_B/\sqrt{W}$ object by object --- to leading Gaussian/log-variance order the object enters only through $K_{\mathrm{eff}}$ --- and normalized by that floor the curves collapse to $\pm5$--$7\%$; under ordered Hadamard the curves scatter by object and do not decay with $W$ --- the object envelope is a bias no amount of averaging removes. The window $W$ carries a bias--variance trade-off (drift curvature versus $\mathrm{CV}_B^2/W$): Sec.~6 gives a general mixing-law upper bound and a matching rate only for the separately specified fixed iid-Gaussian interior experiment.

\begin{figure*}[!t]
\centering
\begin{minipage}[t]{0.30\textwidth}\centering
\includegraphics[width=\linewidth]{fig3a_error_vs_W.png}\\
{\footnotesize (a)}
\end{minipage}\hfill
\begin{minipage}[t]{0.30\textwidth}\centering
\includegraphics[width=\linewidth]{fig3c_collapse.png}\\
{\footnotesize (b)}
\end{minipage}
\caption{Blind gain identifiability (seven frame-matched arms, 2048 frames each; $32\times32$ objects, $K=1024$; medians over 20 seeds $\times$ 10 objects at $\rho_{\mathrm{pair}}=10^{-3}$, $s=0.1$). (a) Blind gain error $\|\hat a-a\|/\|a\|$ versus window $W$ under the ratio AGC (solid, filled) and the log-window implementation (dashed, open). The latter coincides with the Theorem-B estimator only on originally positive, all-one-mask arms; raw signed-Walsh arms are masked diagnostics. Error bars: two-way clustered 95\% CIs of the median best-window error. (b) Errors normalized by the theory floor $\sqrt{\mathrm{CV}_B^2/W}$ collapse across the ten objects. Quantitative values: Sec.~9.2 and Table~III.}
\end{figure*}

\section{Estimation rate, lower bounds, and low-photon robustness}

Section~5 establishes relative-gain consistency under explicit carrier-window conditions. The baseline is the plain windowed log of Theorem~B. For photon counts we instead average $\psi_\alpha(c)=\log(c+\alpha)$ and invert the per-window calibration curve $M_n(\theta)=\sum_kw_{nk}m_\alpha(\Lambda_0e^\theta b_k+d)$ built from \emph{known} carriers $b_k$: an oracle-known-carrier estimator, distinct from the blind high-count estimator. The homogeneous shortcut $m_\alpha^{-1}(\mathrm{mean}_W\psi_\alpha)$ is exact only for a constant carrier. For fixed $\alpha$, the high-count limit has sensitivity tending to one and approaches the ordinary log behavior; no unqualified joint $\alpha\to0$ limit is used. The window trades smoothing bias against count variance.

\begin{namedthm}{Theorem C (oracle-known-carrier finite-window rate)}
Assume conditionally independent Poisson counts across distinct frames, known positive carriers, a nonempty calibration-safe log-gain interval, and $\beta\in(0,1]$. For the endpoint-clamped calibrated inverse,
\[
\mathbb E[(\hat\theta_n-\ell_n)^2\mid\ell,b]
\le(B_n^2+V_n)/\underline\kappa_n^2,
\]
where $B_n$ is the carrier-weighted smoothing bias, $V_n$ the exact conditional soft-log variance, and $\underline\kappa_n$ the window sensitivity. The optimized power law is asserted only when its unconstrained width is feasible and the sensitivity ratios are uniformly controlled. For arbitrary carriers no $\beta>1$ claim is made; sensitivity-balanced interior windows give the partial extension in Appendix~D.4. (Proof: Appendix~D.)
\end{namedthm}

At $d=0$, direct Poisson Le~Cam/Assouad bounds give the matching local order for a centered log-domain loss chosen to match the blind quotient metric, in the locally-flat-gain, feasible-window regime. The oracle experiment itself fixes the absolute scale through known $\Lambda_0$ and carriers. This Poisson statement is proved directly; it is not inherited from the Gaussian theorem below.

\begin{namedthm}{Theorem D (fixed-Gaussian interior minimax rate)}
In the fixed experiment $Y_n=m+f(n/N)+\varepsilon_n$, $\varepsilon_n\stackrel{\mathrm{iid}}{\sim}\mathcal N(0,\sigma^2)$, impose the gauge $N^{-1}\sum_nf(n/N)=0$, the bounded H\"older class $\mathcal H_\beta^0(L,A)$, and pointwise squared loss at an interior grid point. If $h_*=(\sigma^2/(NL^2))^{1/(2\beta+1)}$ is a feasible centered bandwidth and $Lh_*^\beta\le A/4$, then
\[
\inf_{\hat f}\sup_{f\in\mathcal H_\beta^0(L,A)}
\mathbb E(\hat f-f)^2\asymp_\beta
L^{2/(2\beta+1)}(\sigma^2/N)^{2\beta/(2\beta+1)}.
\]
No three-segment finite-$N$ minimax equivalence is claimed outside this interior regime. (Proof: Appendix~D.)
\end{namedthm}

The floor beneath the blind bound is Fisher geometry: in the gauge quotient (global scale unobservable, Sec.~5) the relevant object is the \emph{quotient Fisher information}, which is singular along the exact ambiguity directions generated by Theorem~A. Non-identifiability is not poor conditioning; it is a genuine flat direction of the likelihood --- the exact collision curve $\ell_u=\ell+u$, $T_u=M^{-1}\operatorname{diag}(e^{-u})MT$ made infinitesimal (tall designs need the range condition of Appendix~D.3; constant $u$ is the gauge).

Which design constant controls what is a spikiness question, and three distinct functionals appear: $K_{\mathrm{eff}}$ governs random-carrier excitation (Prop.~1) and also enters the pairwise-Hadamard residual-gain budget (Prop.~3); $K_\infty$ is the Lindeberg parameter gating the Gaussian approximation and windowed-energy concentration, while the log expansion additionally needs the Appendix~C positivity and tail conditions; $K_4$ enters the permutation-whitening concentration guarantee (Sec.~7, via $\min(\epsilon^2 K_4,\epsilon K_\infty)\ge C\log(K/\delta)$). Conflating them mislabels the bottleneck: a design flat in $K_{\mathrm{eff}}$ can still be spiky in $K_\infty$.

Low photon counts do not break the log domain: the Poisson derivative identity $\frac{d}{d\lambda}\mathbb{E}f(C)=\mathbb{E}[f(C{+}1){-}f(C)]$ gives sensitivity $\kappa_\alpha(\lambda)\to 1$ at high $\bar\lambda$ and $\kappa_\alpha(\lambda)\to\lambda\log(1+1/\alpha)$ as $\bar\lambda\to 0$, and the offset $\alpha$ absorbs empty bins so no frame contributes $\log 0$. Through Theorem~C the photon-starved variance factor obeys the general window law of Appendix~D.4 (actual intensities in the numerator, anchor intensities in the sensitivity denominator), simplifying to $\sim1/(W\bar\lambda)$ only under \emph{local gain flatness} --- where it matches the direct Poisson lower bound at rate level, since at $d=0$ the Poisson Fisher information for a known-carrier log-intensity shift equals the expected intensity $\lambda$ and therefore vanishes as $\lambda\to0$; low expected counts cut precision but never make the estimator diverge (Fig.~5). The reference curve plotted there, $\mathrm{mean}_n[1/I_n]$, is an \emph{oracle common-shift diagnostic}, not a nuisance-path quotient Cram\'er--Rao bound.

The caveat, stated honestly: without calibration inversion, a plain windowed average of $\psi_\alpha$ --- the \emph{uncalibrated proxy} --- becomes shrinkage-bias-dominated at $\bar\lambda\lesssim1$: it is pulled toward $\log\alpha$, and its MSE sits \emph{below} the oracle common-shift diagnostic because bias is present (MSE below $1/I_A$ is compatible with shrinkage bias; MSE above it does not certify unbiasedness --- Appendix~D.4). The calibrated estimator of Theorem~C removes this shrinkage and tracks the $1/(W\bar\lambda)$ law without falling below the diagnostic on the tested grid, giving up precisely the bias-bought low-count MSE discount of the proxy (and of the Anscombe transform \cite{Anscombe1948}) --- quantified in Sec.~9 (Fig.~5).

\section{Basis design: randomized orthogonal acquisition (SRHT)}

Must one surrender orthogonal conditioning to buy the anchor's random illumination? No. The randomized Hadamard transform --- an orthogonal Hadamard carrier composed with a random sign flip and a random pixel permutation, the SRHT of the sketching literature --- keeps a full-rank, well-conditioned forward map while randomizing the \emph{temporal handwriting} of the object: one shuffles the deck without bending the cards. We use the \emph{full} $N=K$ randomized transform throughout; ``subsampled'' is reserved for the $N<K$ case, which raises a bias/prior question outside our scope.

The mechanism is exact. Writing the per-coefficient excitation as $Z_g$ under a random diagonal sign $D$, the residual gain--object coupling is governed entirely by the Walsh spectrum of the squared object. With $w = T^2$,
\[\operatorname{Cov}_D(Z_g, Z_h) = \hat{w}(g+h),\]
the Walsh transform of $w$ at the dyadic index sum. Sign randomization alone gives \emph{exact} pairwise whitening iff $\hat w(k)=0$ for every $k\ne0$; it gives $\varepsilon$-whitening iff $\max_{k\ne0}|\hat w(k)|\le\varepsilon S_2$. These are conditions on the object, not the design: a structured $T$ can concentrate $\hat w$ off the origin, and sign flips cannot remove it.

The random permutation supplies the missing genericity, acting as a probabilistic Walsh-flattener: permuting pixels spreads the mass of $\hat{w}$ across the spectrum, driving the off-origin covariance below tolerance $\varepsilon$ with probability $1-\delta$ once the object is spectrally spread enough,
\[\min\!\big(\varepsilon^2 K_4,\ \varepsilon K_\infty\big) \ \geq\ C \log(K/\delta),\]
with $K_4$, $K_\infty$ the Sec.~6 spread functionals --- both diverge as the object's energy delocalizes and collapse when a few pixels dominate; the bound quantifies how much delocalization the permutation needs.

Permutation is a genuine second ingredient. Under permutation alone,
\[\operatorname{Var}_P Z_g = \frac{K S_2 - S_1^2}{K-1} = S_2\,\frac{K - K_{\mathrm{eff}}}{K-1},\]
so, for the physical nonnegative-object regime $K_{\mathrm{eff}}\ge1$, the residual carries an $O(1/K_{\mathrm{eff}})$ \emph{upper} bound on variance but no matching lower bound: a flat object produces zero excitation, so permutation cannot guarantee two-sided control coefficient by coefficient, and the whitening claim is one-sided unless the sign flip and the spectral-spread condition are both in force. A row-order shuffle of the acquisition sequence is a third, weaker operation: it makes the temporal carrier finite-population exchangeable and, under the Serfling range/variance scaling of Appendix~C, concentrates its random-permutation window means; it does not guarantee ($\star$) for every realized order or two-sided per-coefficient excitation. The ablation of Supplementary Material S1 (Fig.~8) probes the acquisition-time operations (none / row-order permutation / sign / both); its permutation arms reorder acquisition time ($P_{\mathrm{row}}$) and do \emph{not} realize the pixel permutation $P_{\mathrm{col}}$ of the whitening analysis above, whose factorized attribution --- including the $P_{\mathrm{col}}$-only and $P_{\mathrm{col}}{+}D$ arms in natural acquisition order --- is reported in Sec.~9.6 and Appendix~E.6.

\begin{nameddef}{Design rule (blind slow gain)}
(i) Do \emph{not} acquire in ordered Hadamard when the nuisance is a blind, slowly drifting gain --- the ordering hands the object a temporal signature (Sec.~4). (ii) Randomize the carrier: apply the SRHT sign flip \emph{and} permutation. (iii) Verify Walsh-flatness of the realized (permuted) $T^2$, or rely on the permutation bound above with $K_4, K_\infty$ estimated from the object class. (iv) Keep a positivity/offset margin so the log or soft-log stage of Sec.~6 stays away from its singularity. (v) Choose the averaging window $W$ from the bias--variance trade of Sec.~6, not from conditioning.
\end{nameddef}

The rule prescribes \emph{both} randomizations because that configuration has a two-sided theorem guarantee under the stated spectral-spread condition; the ablation (Fig.~8) reports the different, empirical fact that either randomization alone already recovered essentially the full advantage in the tested panel. The two statements must not be merged: the theorem guarantees whitening under its hypotheses, the ablation shows that any tested way of breaking the chronology sufficed here; we do not promote the panel result to a universal claim.

The synthesis: orthogonality and randomized temporal homogenization are compatible. SRHT preserves the $B_L=1$ leverage of a full orthogonal inversion while, under Walsh-flatness or the stated permutation concentration event, controlling cross-row covariance over the pixel-randomization ensemble. A separate $P_{\mathrm{row}}$ randomization gives finite-population exchangeability of acquisition order. Neither statement makes a conditioned fixed draw stationary or mixing; the realized window criterion remains a check.

(Proof of the Walsh identity and permutation bound: Appendix~E.)

\section{From identifiability to images: the reconstruction bridge and phase diagram}

Sections~4--7 answer the statistician's question --- \emph{is the gain identifiable, and at what rate?} --- but the practitioner asks: \emph{how many gray levels do I lose?} This section supplies an exact finite-$N$ second-moment bridge under its stated exogeneity hypotheses. Physical Poisson and signed-coefficient implementations enter through separately derived covariance models.

The corrected bucket record is $z_n=(1+\delta_n)b_n+\xi_n$ with $b=AT$, $\delta_n$ the \emph{residual} multiplicative gain after whatever blind correction the pipeline applied, and $\hat T=Lz$ any linear reconstruction.

\begin{namedthm}{Theorem 1 (master finite-noise relMSE identity)}
Assume --- as part of the hypothesis, not as background convention --- that (H1) $\mathbb{E}\,\xi=0$; (H2) the residual gain $\delta$ and the additive noise $\xi$ are uncorrelated, $C_{\delta\xi}:=\operatorname{Cov}(\delta,\xi)=0$; (H3) the gain estimate $\hat a$ is \emph{exogenous} with respect to the reconstructed record --- computed independently of the noise realization $\xi$ (injected or synthetic residuals, an independent calibration record, or a cross-fitted estimate from held-out frames), so that (H1)--(H2) hold conditionally on the record being reconstructed; and (H4) the second moments $m_\delta$, $V_\delta$, $\Sigma_\xi$ are finite. Then, without Gaussian or small-noise approximation,
\begin{align*}
\frac{\mathbb{E}\|\hat T-T\|^2}{S_2}\;=\;&\underbrace{\frac{\|(LA-I)T+L\,\mathrm{diag}(b)\,m_\delta\|^2}{S_2}}_{\text{bias}}\\
&+\;\underbrace{\frac{\operatorname{tr}\!\big(L\,\mathrm{diag}(b)\,V_\delta\,\mathrm{diag}(b)\,L^\top\big)}{S_2}}_{\text{residual gain}}\\
&+\;\underbrace{\frac{\operatorname{tr}\!\big(L\,\Sigma_\xi L^\top\big)}{S_2}}_{\text{additive noise}}.
\end{align*}
If the residual gains are independent with $\mathrm{Var}\,\delta_n=v$, the gain term collapses to $v\,B_L$ with the \textbf{leverage} $B_L=(1/S_2)\sum_n b_n^2\|Le_n\|^2$; coherent (smooth) residuals require the full matrix form. (Proof: Appendix~F.)
\end{namedthm}

The hypotheses are not decorative: a gain estimate computed from the \emph{same} noisy record carrying $\xi$ --- the generic blind situation --- generically violates (H1)--(H3), adding an $L\,m_\xi$ bias term and a $C_{\delta\xi}$ cross-term (Remark~F.1.1). The identity is accordingly validated in the injected-residual (exogenous) regime (E4, Sec.~9) and applies to cross-fitted estimates restoring (H1)--(H2) by construction; for same-record blind pipelines it is exact only up to those two terms.

For physical nonnegative measurement rows, the Poisson plug-in $\Sigma_\xi=\mathrm{diag}(\tau_{G,n}^2+b_n^2/\Phi_n)$ is exact at all counts after the post-calibration substitution $\hat a_n\to a_n$; before it both variances carry $(a_n/\hat a_n)^2$. A signed transform coefficient $b_n=(UT)_n$ is not a Poisson intensity: complementary masks or a positive offset require their own photon allocation and differencing covariance. The identity uses raw error $\mathbb E\|\hat T-T\|^2/S_2$. Fig.~9 is therefore regenerated with raw total MSE for every theorem-validation arm; scale-aligned MSE remains a separately labeled image-quality output. Specializing $L$:

\begin{itemize}[leftmargin=*,itemsep=1pt,topsep=2pt]
\item \textbf{Orthogonal / SRHT full inversion:} $B_L=1$, so independent residual gain passes through at unit leverage. For signed-coefficient additive covariance $\Sigma_c$, $\mathrm{relMSE}=v+\operatorname{tr}(\Sigma_c)/S_2$; no direct signed-Poisson term is inserted.
\item \textbf{Pairwise (differenced) Hadamard:} gain term $R_{\mathrm{pair}} = (K_{\mathrm{eff}}/4)\,D_H(T)\,\mathbb{E}[\Delta^2]$ --- the second moment $\mathbb{E}\Delta^2$, not $\mathrm{Var}(\Delta)$, unless $\mathbb{E}\Delta=0$ is separately assumed (a deterministic mismatch $\Delta\equiv\epsilon$ has zero variance but nonzero risk; Appendix~F.3.2) --- with $D_H\approx 1$ for non-DC coefficients and $\mathbb{E}[\Delta^2]\approx 2s^2\,r(\rho_{\mathrm{pair}})$ the adjacent-pair increment for OU drift ($r(\rho)=1-e^{-\rho}$; $r(\rho)\approx\rho$ for small $\rho$): differencing converts \emph{slowness} directly into suppression.
\item \textbf{Random / DGI:} $\mathrm{relMSE}_{\mathrm{DGI}} = C_0^{\mathrm{floor}}/N + (1+C_0^{\mathrm{lev}})v/N + $ noise terms. Three constants of the \emph{same} pipeline must be kept apart (Appendix~F.3.3): the fixed-design sampling floor $C_0^{\mathrm{floor}}$ (raw $\approx10^3$ here, aligned $\approx7\times10^2$ --- a factor-1.54 metric difference), the drift leverage $\Lambda_{\mathrm{lev}}=NB_L=1+C_0^{\mathrm{lev}}$, and the population moment constant $K+\beta_4-2$. Mean subtraction removes the $K_{\mathrm{eff}}K(\mu/\sigma)^2$ background penalty from the floor but \textbf{not} from the leverage --- drift modulates the large DC background coherently --- so for mean-subtracted pipelines $C_0^{\mathrm{lev}}$ is the nonnegative-background value of (F.10), roughly $10^3\times$ the floor. There is \textbf{no universal $C_0$}: each constant depends on pattern statistics, background, normalization, \emph{and metric} --- for our SCGI-style testbed \cite{PengChen2024APL} the measured values are used throughout Fig.~9 and Fig.~4.
\end{itemize}

In the square unit-leverage comparisons studied here, orthogonal inversion is the reference conduit: $B_L=1$ and there is no reconstruction floor. This is not a universal optimum over redundant designs --- repeated or averaged measurements can have $B_L<1$. The randomized arms win here primarily because the stationarity anchor lets the residual error itself shrink, where ordered chronology offers no comparable blind handle.

These closed forms compete, producing a finite-$N$ flip boundary.

\begin{namedthm}{Prop 3 (finite-$N$ flip boundary; fixed-point form)}
All risks are the \emph{raw} relative MSEs of Theorem~1. Write the random-arm risk as $R_{\mathrm{rand}}=C_0^{\mathrm{floor}}/N+\Gamma_{\mathrm{blind}}+R_{\mathrm{rand,noise}}$, where $\Gamma_{\mathrm{blind}}$ is the full residual-gain contribution. For an exogenous, centered, frame-uncorrelated residual only, $\Gamma_{\mathrm{blind}}=(1+C_0^{\mathrm{lev}})q_{\mathrm{blind}}/N$. A same-record blind estimator requires the full projected mean, covariance, and residual--noise cross terms and has no asserted scalar $v_{\mathrm{blind}}$ representation. In the scalar special case, $\rho^{*}=\inf\{\rho_{\mathrm{pair}}\ge0: R_{\mathrm{pair}}\ge R_{\mathrm{rand}}\}$ satisfies
\[
\begin{aligned}
1-e^{-\rho^{*}}=\frac{2}{K_{\mathrm{eff}}\,D_H\,s^{2}}\big[&C_0^{\mathrm{floor}}/N\\
&+\Gamma_{\mathrm{blind}}(\rho^*,N)\\
&+R_{\mathrm{DGI,noise}}-R_{\mathrm{pair,noise}}\big],
\end{aligned}
\]
Because the right side moves with $\rho^*$, this is not an explicit solution: existence needs a crossing and uniqueness a single-crossing condition. Only a constant bracket $Q_0$ gives $\rho^*=-\log(1-Q_0)$; $Q_0\ge1$ gives no finite crossing. Its gain-known leading heuristic is $\rho^{*}\approx2C_0^{\mathrm{floor}}/(NK_{\mathrm{eff}}s^2)$. (Proof: Appendix~F.)
\end{namedthm}

\emph{Remark:} $\rho_{\mathrm{pair}}$ (adjacent-pair decorrelation, this section) and $\rho_{\mathrm{bw}}$ (normalized bandwidth, Sec.~4) both grow with drift speed but are not equal; every Fig.~4 axis declares its convention (Appendix~F.4).

The injected-residual bridge is tested in raw MSE in Supplementary Material S2 (Fig.~9); because the protocol contains no photon noise, it isolates residual-gain propagation. Figure~4 reports the equal-frame competition and the no-free-parameter gain-known skeleton. The same-record blind no-flip result is an empirical observation accompanied by a conditional noiseless projected-covariance diagnostic, not a scalar theorem validation. Three operational regimes still organize the chart: slow drift favors pair differencing, intermediate drift favors randomized chronology, and fast drift defeats blind separation unless an external gain reference is supplied.

The measured constants and conventions behind Fig.~4 are collected in Table~I.

\begin{table*}[!t]
\centering
\caption{Headline measured constants behind Fig.~4 (secondary constants and grid conventions: Supplementary Material Table~II). Every constant is pipeline-specific, and floor, leverage, and moment constants of the \emph{same} pipeline are different numbers (Appendix~F.3.3). Confidence intervals: two-way (seeds-and-objects) crossed cluster bootstrap percentile intervals over the 20-seed $\times$ 10-object replication, slope fits restricted to $W\le16$; an object$\times$seed bootstrap treating the 200 cells as independent understates uncertainty; no qualitative conclusion changes.}
\label{tab:constants}
\begin{tabular}{@{}p{0.24\textwidth}p{0.50\textwidth}p{0.20\textwidth}@{}}
\toprule
Quantity & Value in this pipeline & Uncertainty \\
\midrule
Drift leverage $\Lambda_{\mathrm{lev}}=NB_L=1+C_0^{\mathrm{lev}}$, distinct from the Prop.~3 raw floor & $\Lambda_{\mathrm{lev}} \in (7.527\text{--}7.544)\times10^{5}$ across $N$ & exact fixed-operator column-norm calculation, then mean over 10 objects; no residual probe \\
\addlinespace
$D_H(T)$ (pairwise-Hadamard factor) & measured over the ten-object panel: min/median/max $=0.971/0.988/0.999$ (per object via the analytic identity $D_H(T)=\frac{1}{K}\sum_k(1-x_k^2)^2$, $x_k=c_k/S_1$) & deterministic per object (identity evaluated on the panel); no resampling needed \\
\addlinespace
Same-record blind residual & report raw $q_\delta$, residual mean, lag correlation, and projected covariance; scale-aligned gain error is secondary & conditional noiseless proxy, not a scalar theorem input \\
\addlinespace
Prop.~3 boundary, raw metric on both sides & true-$S_1$: $40/40$, median/max factors $\mathbf{1.020}/1.037$, slope $-2.12044$; production median normalization: $40/40$, $1.044/1.200$, slope $-2.00656$; prediction $-2.11371$ & same-estimator validation versus production robustness, identical OU traces; clean-commit authoritative R15 result \\
\bottomrule
\end{tabular}
\end{table*}

\begin{figure*}[!t]
\centering
\includegraphics[width=0.56\textwidth]{fig5_flip_boundary.png}
\caption{Empirical winner map, SRHT crossover, and the raw-MSE Prop.~3 skeleton test at equal frame budget ($N=2048$, $K=1024$, OU drift). (a) Winner per $(\rho_{\mathrm{pair}},\sigma_a)$ cell among candidates passing the above-floor gate (grey: floor; outline: the single non-SRHT winner). (b) Descriptive SRHT-paired-versus-ordered-Hadamard crossover fits, a comparison Prop.~3 does not rank. (c) No-free-parameter gain-known skeleton: prediction and factor-2 band against the theorem-faithful true-$S_1$ arm (filled circles) and production median normalization (open squares) on identical traces; triangles: no crossing in range. True-$S_1$ median/max factors are $1.020/1.037$; production factors $1.044/1.200$.}
\end{figure*}

\section{Numerical experiments}

The theory of Secs.~4--8 makes quantitative predictions --- a stationarity dichotomy, a leading Gaussian/log-variance $W^{-1/2}$ law whose tested-panel object dependence is summarized by $K_{\mathrm{eff}}$, an exact relMSE identity, a flip boundary, and a photon-limited Fisher scaling --- and this section tests each. A closing subsection integrates five later verification experiments: the tall-design rank threshold (Fig.~6), whitening power (Supplementary Material S1), the coherent-residual validation of Theorem~1 (S2), the $C_0$ audit, and oracle-to-blind baselines (S3).

\noindent\textbf{Protocol.} All results are simulations; no hardware data enter. Comparisons use matched frames and flux. Each seed shares one OU log-gain trace across arms. The H\"older model is the analysis class, whereas OU is the simulation protocol and is not pathwise Lipschitz. Ten fixed objects span $K_{\mathrm{eff}}=59.8$--$864.3$; the base seed is $20240708$. Blind gain error is scale-aligned after unit-mean normalization; mask contrast is GT-mask CNR. Theorem~1 and Prop.~3 validation use raw reconstruction MSE, while scale-aligned reconstruction MSE is retained only as a separately labeled image-quality metric. Full enumerables: Supplementary Material S5.

\noindent\textbf{Carrier stationarity (Fig.~2).} This experiment applies a \emph{variance-stationarity diagnostic} to the noiseless carrier $\{B_n\}$ before any gain is applied (8192 frames per arm, $64\times64$ objects). The Brown--Forsythe test checks chunk-variance equality: it screens the stronger strict-stationarity sufficient model used in (B2), not the windowed log-mean condition ($\star$) itself and not a necessary condition for ($\star$). The test at $p<10^{-3}$ rejects 7 of 10 objects for ordered Hadamard, with $p$ reaching $10^{-20}$, against 0 of 10 for both i.i.d.\ random ensembles (uncorrected per-test $p$-values; the strongest rejections clear a Bonferroni-adjusted threshold by many orders of magnitude). The permuted arms --- randomly permuted paired Hadamard ($P_{\mathrm{row}}$) and SRHT ($P_{\mathrm{row}}{+}D$, Appendix~E.6) --- each reject 2 of 10: a \emph{fixed} reordering draw does not guarantee temporal homogeneity for every object, since the randomization guarantees of Secs.~5 and 7 are probabilistic over draws --- this residue is the theory's own prediction, not an anomaly. The pixel-level randomization of the SRHT analysis ($P_{\mathrm{col}}$, $D$) is probed by the factorized ablation of Sec.~9.6. The Parseval invariant closes the loop: every exhaustively paired arm has carrier $\mathrm{CV}_B = 1/\sqrt{K_{\mathrm{eff}}}$ to about $10^{-6}$ relative accuracy, and the random arms match $(\sigma_I/\mu_I)/\sqrt{K_{\mathrm{eff}}}$ within $\pm2\%$ --- the ordered pathology lives entirely in the temporal variance structure, not in the marginal spread (descriptive counts over the ten-object panel).

\noindent\textbf{Blind gain identifiability (Fig.~3).} This is the core evidence for Theorem~B, run over nine arms including the raw, non-paired Hadamard arms that feed bare signed Walsh coefficients to the estimator (bucket convention (i), diagnostic-only); per-arm errors and slopes with two-way CIs: Supplementary Material Table~III. Raw arms are scored at the same 2048-frame budget; the single-pass 1024-frame diagnostics are indistinguishable (Table~III), so the failure is chronology, not budget. In Fig.~3 the chronology failure of Theorem~A is made visible: under the ratio AGC the raw ordered arm's blind gain error is $0.94$--$0.99$ across all averaging windows (median best-window $0.899$) and does not decay; shuffling the raw frames improves this to $0.794$ --- permutation helps even here --- but the arm remains far above the random and paired arms at $0.007$--$0.017$, because without pairing the near-zero-mean carrier leaves the AGC ill-posed. The log-window implementation coincides with the Theorem~B estimator only on arms whose original analyzed carrier is strictly positive (so the positivity mask is all one); on raw signed arms it is an endogenous masked diagnostic, not a theorem test. It agrees with the ratio AGC to about $2\%$ on all positive stationary-carrier arms with the same qualitative ordering; on the raw diagnostics, ordered still cannot fall below $0.71$ (non-decaying) while shuffled improves to $0.136$, still at least $12\times$ above every randomized arm. Within the variance-obeying arms, the slow-drift variance-segment slopes of $\log(\mathrm{err})$ against $\log W$ all sit inside the $W^{-1/2}$ band ($-0.39$ to $-0.61$), while naturally ordered paired Hadamard is shallow at $-0.24$: differencing cancels the drift within pairs but leaves the ordered envelope. Normalizing by the theory floor $\sqrt{\mathrm{CV}_B^2/W}$ collapses the panel: the paired arms sit at $0.40$--$0.41$ and the random arms at $0.89$--$0.90$, each bundle within $\pm5$--$7\%$ across all ten objects, and the best-window error correlates with $K_{\mathrm{eff}}$ at $-0.86$ to $-0.90$ (descriptive) --- consistent, within this tested panel, with the object entering through $K_{\mathrm{eff}}$ alone, as Proposition~1 predicts to leading Gaussian/log-variance order. The shuffled-paired arm performing like the random arms bears out the permutation analysis of Sec.~7: pairing cancels the slow drift within pairs, and the permutation restores the stationarity that the ordering destroyed.

The 20-seed uncertainty replication (31{,}600 rows) confirms these slopes under the fixed $W\le16$ variance-segment convention: the randomized arms' median two-way slopes lie between $-0.39$ and $-0.62$, the naturally ordered paired arm's ratio-AGC slope CI straddles zero (no detectable decay), and the raw ordered arm stays essentially flat at $0.899$ (per-arm values and CIs: Table~III). The two-way crossed cluster bootstrap widens naive object$\times$seed intervals by $3$--$5\times$ without crossing any qualitative boundary (rationale: Table~I caption). One nuance: under the log-domain estimator the naturally ordered paired arm shows a shallow decay ($-0.222$ $[-0.306,-0.054]$) where the ratio AGC shows none; both remain far outside the $W^{-1/2}$ band, so the chronology conclusion stands (``no detectable decay'' attaches to the ratio AGC). Best-window blind gain errors at this operating point are in Table~I; the exploratory slopes and $R^2$ above used an adaptive variance segment, and the two conventions agree on every qualitative conclusion.

\noindent\textbf{The reconstruction bridge (Supplementary Material S2, Fig.~9).} The R15 bridge uses one metric and one quantity in every arm: total \emph{raw} relative MSE versus the fixed-design prediction $R_{\mathrm{floor}}^{\mathrm{raw}}+B_Lv$. The balanced 8100-row protocol contains three reconstructors, $N\in\{1024,2048,4096\}$, six $v$ levels (including the $v=0$ baseline), ten objects, and fifteen seeds. The deterministic inverses remain flat in $N$ (per-$v$ slopes $[-0.117,0.112]$ orthogonal and $[-0.014,0.009]$ SRHT), while random/DGI follows $1/N$ (slopes $[-1.005,-0.972]$). Across the $(N,v)$ aggregates, measured/theory ratios lie in $[0.913,1.163]$ for the orthogonal inverse, $[0.992,1.021]$ for SRHT, and $[0.982,1.010]$ for random/DGI; exact fixed-operator leverage gives $NB_L=(7.527\text{--}7.544)\times10^5$. The former DC-object discrepancies were created by scale alignment and disappear from the raw theorem panel. The per-realization identity $R_{\mathrm{raw}}=R_{\mathrm{floor}}+R_{\mathrm{gain}}+R_{\mathrm{cross}}$ closes to $2.28\times10^{-7}$ absolute error. The clean-commit r3 output is the submission-authoritative bridge dataset.

\noindent\textbf{Phase diagram and ablation (Fig.~4; ablation panels in Supplementary Material S1, Fig.~8).} The phase diagram sweeps $\rho\in[10^{-3},10]$ and $\sigma\in[0.05,0.5]$ at equal frames. Of 29 above-floor cells, SRHT-paired acquisition wins 28; one goes to randomly ordered pairwise Hadamard with SCGI-proxy correction. The plotted two-parameter crossover fits summarize SRHT-paired versus ordered Hadamard and are not a Prop.~3 test. The no-free-parameter raw-MSE test is separate and now has two explicit estimators on identical OU traces. The theorem-faithful true-$S_1$ arm observes 40/40 crossings at $\sigma_a\ge0.1$, with median/max agreement factors $1.020/1.037$ and pooled slope $-2.12044$ versus the predicted $-2.11371$. The production median-normalized arm also observes 40/40, with factors $1.044/1.200$ and slope $-2.00656$; it is governed by (F.7a) and is a robustness check, not the same-estimator validation. The previous factor $1.54$ was a raw/aligned mismatch. The blind random arm never overtakes ordered pairwise Hadamard in the 100 scanned cells. Because its estimator uses the same noiseless record, the R15 analysis reports raw $q_\delta$ and the exact projected mean/covariance decomposition and retains the scalar no-flip curve only as a conditional proxy; no detector-noise cross term is tested. The intermediate-drift winner is SRHT-paired acquisition. The ablation at $\rho=10^{-3}$ gives $+5.45$ dB for SRHT, $+5.47$ dB for pixel signs alone, and $+5.39$ dB for acquisition-order permutation alone (descriptive, nonadditive). At $\rho\ge1$ all blind variants approach the noise floor within $\pm0.09$ dB; this is not a direct low-pass-threshold test unless $\rho_{\mathrm{bw}}$ and $\rho_{\mathrm{pair}}$ are controlled separately.

\noindent\textbf{Low-photon robustness (Fig.~5).} This experiment tests the soft-log analysis of Sec.~6 under Poisson counting, with \emph{five} estimator arms on the same paired Poisson draw at each budget: naive clipped log, Anscombe, the uncalibrated soft-log proxy, the mean-Poisson calibrated proxy $m_\alpha^{-1}(\mathrm{mean}_W\,\psi_\alpha)$, and the oracle-known-carrier calibrated soft-log of Theorem~C (per-window carrier-calibration root, carriers $b_k$ from the simulation truth --- a flat-field benchmark; requested $W=64$, effective interior width 65; truncated edge windows; Appendix~D.4); each point aggregates 50 object-by-seed cells, error bars the SD over those cells. Every ratio below is a log-domain (centered-log-gain) mean-MSE ratio, numerator over denominator as named --- the metric comparable to the oracle common-shift diagnostic $\mathrm{mean}_n[1/I_n]$ of Appendix~D.4 ($d=0$). At mean counts $\bar\lambda\le1$ the naive clipped log is $24.9$--$28.7\times$ worse than the uncalibrated proxy, $16.3$--$21.0\times$ worse than Anscombe, and $1.8$--$11.4\times$ worse than the calibrated estimator (all naive-over-method ratios). The two calibrated arms track the $1/(W\bar\lambda)$ law across the variance regime --- slopes $-1.000$ (oracle-carrier) and $-0.99$ (mean-Poisson proxy, overlapping the oracle arm to $0.4\%$ at this protocol's small carrier variation) over $\bar\lambda\in[2,32]$, against $-0.920$/$-0.732$ for the uncalibrated proxy (its $[1,16]$ fit is contaminated by its own bias transition). At budgets $\bar\lambda\le2$ the uncalibrated proxy and Anscombe fall \emph{below} the common-shift diagnostic --- shrinkage bias, not super-efficiency (MSE above it likewise does not certify unbiasedness; Appendix~D.4) --- whereas the calibrated estimator stays above it at every tested budget, within $2$--$18\%$ at $\bar\lambda\le1$ and $14$--$31\%$ through $[2,32]$. The proxy's apparent $1.19$--$1.53\times$ advantage over Anscombe at these budgets is that shrinkage bias, not efficiency; and the calibrated estimator loses to Anscombe in mean MSE at $\bar\lambda\le1$ (Anscombe $1.8$--$8.9\times$ better) --- a bias-bought advantage (Sec.~6 trade-off) --- overtaking it only for $\bar\lambda\ge32$, where the pairwise ratios converge to $1.00$ ($\pm4\%$). At the high-photon end the floor is drift-limited, not photon-limited: $2.434\times10^{-4}$ ($\rho=10^{-3}$) $\to$ $1.411\times10^{-4}$ ($\rho=10^{-4}$) for the proxy arms, $2.168\times10^{-4}\to1.125\times10^{-4}$ for the oracle-carrier estimator --- $10.92\%$ lower at $\rho=10^{-3}$ and $20.28\%$ at $\rho=10^{-4}$ (two distinct margins), since it deconvolves the known carrier (floor probe: Supplementary Material S4, Fig.~12).

\begin{figure}[!t]
\centering
\includegraphics[width=0.46\columnwidth]{fig7_logMSE_vs_photon.png}
\caption{Low-photon robustness under Poisson counting, on identical random-basis carriers and shared gain traces (soft-log $\alpha=0.5$, requested $W=64$ / effective interior width 65, $\rho_{\mathrm{pair}}=10^{-3}$, $s=0.1$; truncated edge windows). Five methods --- naive clipped log, Anscombe, uncalibrated soft-log proxy, mean-Poisson calibrated proxy, oracle-known-carrier calibrated soft-log (Theorem~C) --- on the same paired Poisson draws; error bars: SD over the 50 object-by-seed cells. Centered-log-gain MSE ($d=0$) versus $\bar\lambda$ against the \emph{oracle common-shift diagnostic} $\mathrm{mean}_n[1/I_n]$ (dashed; dotted: nominal $1/(W\bar\lambda)$; shaded: $\bar\lambda<1$) --- not a nuisance-path quotient CRB. The calibrated arms (diamonds: oracle-carrier root; the mean-Poisson proxy overlaps it) track the $1/(W\bar\lambda)$ law and stay above the diagnostic; the uncalibrated proxy and Anscombe buy low-count MSE with shrinkage bias (below the diagnostic at $\bar\lambda\le2$). Ratios, slopes, floor: Sec.~9.5.}
\end{figure}

\noindent\textbf{Threshold, whitening-power, and accountability verifications (Fig.~6 and Supplementary Material S1--S4).} Five further experiments close the remaining theory--evidence gaps; unless a bootstrap interval is stated, every number below is a descriptive result of its grid.

\emph{Tall-design threshold scan (Fig.~6).} A direct test of the local component of Theorem~A$'$ at two object sizes ($K=64$: 20 seeds per cell, 1300 cells; $K=128$: 30 seeds per cell, 1950 cells), over five low-pass drift dimensions $p\in\{3,5,9,17,33\}$ at margins $N-K-p\in[-8,+16]$. At both sizes the local rank test on the collision Jacobian transitions from $0\%$ to $100\%$ pass \emph{exactly} at margin $N-K-p=-1$ for every $p$ (3250 cells, no exceptions) --- matching the proved a.e.\ rank formula $\operatorname{rank}D\Phi_M=\min\{N,K+p-1\}$ of Theorem~A$'$(i); below the wall the observed rank deficits correspond, by the same theorem, to genuine positive-dimensional alias families, not numerical degeneracy. Being a rank statement, the scan verifies the local component only, not the global sufficiency thresholds of A$'$(ii)--(iii) or their conjectured sharpness. The companion solver panel (blind alternating recovery, 8 restarts $\times$ 150 iterations) starts at the same wall but climbs slowly --- about $2\%$ success at margin $0$ and $\sim50\%$ at margin $+16$ for $K=64$ ($\sim41\%$ for $K=128$) --- an \emph{algorithmic-recovery} observation, not a uniqueness test: uniqueness and recovery separate, as Sec.~2's third notion anticipates.

\emph{Multi-permutation whitening power and factorized randomization ablation (Supplementary Material S1, Fig.~7).} Two randomization semantics must be kept separate (Appendix~E.6): the pixel-level randomization of the SRHT theory ($P_{\mathrm{col}}$, $D$; acquisition-order-invariant by Lemma~E.1) and the acquisition-order reordering $P_{\mathrm{row}}$, which is what the Brown--Forsythe screen probes. The factorized ablation attributes the whitening-power headline: either permutation alone substantially reduces the rejection rate --- from $70\%$ (ordered) to $5.6\%$ ($P_{\mathrm{col}}$) and $12.2\%$ ($P_{\mathrm{row}}$), the theory's $P_{\mathrm{col}}{+}D$ to $11.9\%$ in natural order --- a large but incomplete suppression (all three stay well above the $\sim0.1\%$ nominal rate of the $p<10^{-3}$ test), so Appendix~E's whitening claims do not depend on the row-order confound; pixel signs alone barely move it ($43.4\%$). The per-draw Walsh certificate fails (Spearman $0.092$, $p=0.10$): an honest null, consistent with the probabilistic nature of permutation whitening (Sec.~7).

\emph{Coherent-residual validation of Theorem~1 (E4, Supplementary Material S2, Fig.~10).} Under five injected residual-covariance structures (150 cells, raw identity-verification metric), the matrix law of Theorem~1 tracks the measured relMSE within $4.2\%$, while the scalar $v\,B_L$ law --- exact for the orthogonal and SRHT inversions --- fails by up to $3.65\times$ for random/DGI under low-pass-coherent residuals: a direct verification of the coherent clause. E4 uses \emph{injected} residuals; it does not probe the same-record $\delta$--$\xi$ dependence flagged after Theorem~1.

\emph{$C_0$ audit and oracle-to-blind baselines (Supplementary Material S3, Fig.~11).} The measured one-sample constant agrees with the Appendix-F moment formula to $0.35\%$ mean ($2.01\%$ worst) relative error over four ensembles $\times$ five objects, the ensemble spread confirming that $C_0$ is pipeline-specific; and on four held-out objects the ordering static $\ge$ oracle $\ge$ blind-AGC $\ge$ no-correction holds in CNR for 4 of 4, with an oracle$-$blind gap of only $0.035$--$0.147$: the blind estimator concedes little to a gain oracle; forgoing correction costs the most.

\begin{figure}[!t]
\centering
\includegraphics[width=0.55\columnwidth]{fig8_threshold.png}
\caption{Tall-design threshold scan (local-rank test of Theorem~A$'$(i) plus an algorithmic-recovery probe) at $K=64$ (20 seeds/cell, 1300 cells) and $K=128$ (30 seeds/cell, 1950 cells), $p\in\{3,5,9,17,33\}$, margins $N-K-p\in[-8,+16]$, all $p$ pooled. (a) The rank test transitions $0\%\to100\%$ \emph{exactly} at margin $-1$ at both sizes (3250 cells, no exceptions), verifying $N=K+p-1$. (b) The blind alternating solver starts at the same wall but climbs slowly --- algorithmic recovery, not uniqueness; rates in Sec.~9.6.}
\end{figure}

\section{Discussion: significance, scope, and outlook}

The results of Sections~2--9 are pre-acquisition design tools, evaluated on the illumination basis and the expected drift \emph{before} photons are collected. Four deliverables follow, then limits and the invited experimental program.

\noindent\textbf{A go/no-go identifiability predictor.} The tall-design thresholds of Theorem~A$'$ convert ``will this acquisition self-calibrate?'' into arithmetic. For $K\ge2$ and almost every design--object pair with nonvanishing carrier, ordinary local identifiability holds iff $N\ge K+p-1$; generic exact sufficiency begins at $N\ge K+p$, and uniform sufficiency at $N\ge 2K+p-1$, for any fixed drift class $S\ni\mathbf 1$ including the physical low-pass family (sharpness of the two global thresholds remains a labeled prediction). The $K=1$ stratum is globally identifiable on the nonvanishing-carrier set at every admissible $N$. Through $p\approx\rho_{\mathrm{bw}}N$, slow drift needs only modest oversampling, while $\rho_{\mathrm{bw}}\to1$ defeats blind algebraic separation and forces an external reference, an additional prior, or another non-algebraic information source. Theorem~B is one such statistical route only when a coherence window satisfying (B3)--(B4) still exists; it also degrades once gain decorrelates within that window (the $\rho\ge1$ convergence in Fig.~8).

\noindent\textbf{An SRHT design rule.} When blind slow gain is the operative nuisance, do not run ordered Hadamard: its chronology hands the object a temporal signature indistinguishable from drift. Randomize the carrier per the Sec.~7 rule, keeping the two permutation semantics distinct (Appendix~E.6): $P_{\mathrm{col}}$ enters the covariance-whitening guarantee, $P_{\mathrm{row}}$ the exchangeability anchor of Sec.~5; empirically either permutation alone substantially reduces the non-stationarity screen's rejections, signs alone barely (Sec.~9.6). Retain a positivity/offset margin and choose $W$ from the bias--variance tradeoff: statistical self-calibration without surrendering orthogonal conditioning.

\noindent\textbf{A design map exposed to falsification.} Theorem~1 and Prop.~3 organize raw MSE over basis, drift, and noise, with $\rho_{\mathrm{pair}}$ distinct from $\rho_{\mathrm{bw}}$. The true-$S_1$ same-estimator arm validates the gain-known skeleton at median factor $1.020$; the production median normalization gives $1.044$ and needs its additional gauge and cross terms. The earlier factor $1.54$ was a raw/aligned mismatch. For same-record blind DGI, the full projected covariance is the rigorous empirical object; the scalar no-flip curve is only a conditional proxy consistent with the 100/100 observed no-flip cells. Floor and leverage remain pipeline- and metric-specific.

\noindent\textbf{A cross-domain principle.} The exportable statement is basis-agnostic: \emph{to self-calibrate a multiplicative nuisance, design the carrier so its local averages are time-homogeneous and concentrate in the nuisance coordinate.} Blind gain/phase calibration, flat-field correction, coded-aperture and computational microscopy, and single-pixel spectroscopy share this bilinear structure --- analogues and motivation, not proven domains; each needs its own operator analysis of the diagonal-row map \cite{LiLeeBresler1501.06120,KechKrahmer1603.07316}.

\noindent\textbf{Scope and limits.} The claims are theory and simulation only, with no hardware validation. Uniqueness near $N=K+p$ is an identifiability statement, not a conditioning guarantee (stability is a separate random-matrix question), and sharpness of the two global thresholds remains a labeled prediction (Sec.~4). Random square unconstrained designs are \emph{not} algebraically identifiable (Theorem~A). The soft-log avoids the log-zero singularity but cannot manufacture information: for a Poisson log-intensity parameter, Fisher information equals the expected intensity and vanishes as $\lambda\to0$ --- an information limit, not an estimator defect.

\noindent\textbf{Experimental feasibility.} The theory issues sharp, falsifiable bench tests under programmable drift: the gain-known flip point $\rho^*$, the conditional no-flip check for mean-subtracted same-record blind DGI, the SRHT-versus-ordered crossover, and the Fig.~3 blind gain-error separation; hardware tests are left to future work with an optical-imaging group.

In conclusion: the square-design obstruction and proved tall-design thresholds (Theorems A, A$'$), triangular-array relative-gain consistency (Theorem B), the fixed-Gaussian interior minimax result and separate oracle-Poisson theorem (Theorems C--D), randomized-Hadamard covariance design, and the raw finite-noise bridge make gain observable not by measuring it directly, but by making the object stop pretending to be time.

\section*{Supplementary Material}

The supplementary document contains: S1 --- permutation-whitening power and the factorized randomization ablation; S2 --- the reconstruction-bridge panel (Fig.~9) and the coherent-residual (E4) validation of Theorem~1; S3 --- the $C_0$ audit and oracle/metered/blind baselines; S4 --- extended flip-boundary detail, the drift-floor probe, and replication tables; S5 --- the shared simulation protocol; Appendices A--F --- the complete proofs. Pointers ``Appendix~X'' refer to the supplementary document; its reference list shares this one's numbering.

\section*{Code and Data Availability}

The project repository contains the experiment runners, committed result CSVs, and \texttt{make\_pub\_figures.py}: \url{https://github.com/ccyyyYyzz/GI_a1} (branch \texttt{scgi-ceiling-diagnostic-r1}). The clean-commit R15 Prop.~3 and raw-bridge outputs used here are committed at \path{results/prop3_estimator_contract_r15/} and \path{results/paper_fig4_bridge_r3_raw/}; their manifests identify source commit \texttt{e4d4996} and classify both as submission-authoritative.
